\documentclass[11pt,a4paper]{article}

% --- Codificación y tipografía ---
\usepackage[utf8]{inputenc}
\usepackage[T1]{fontenc}
\usepackage[spanish,es-noshorthands]{babel}
% --- Fuente más legible: Lato (texto sans-serif) + Inconsolata (monoespaciada) ---
\usepackage[default]{lato}
\usepackage[scaled=0.95]{inconsolata}

% --- Diseño de página, tablas y gráficos ---
\usepackage[margin=2.3cm]{geometry}
\usepackage{array}
\usepackage{booktabs}
\usepackage[table]{xcolor}
\usepackage{titlesec}
\usepackage{enumitem}
\usepackage{tcolorbox}
\usepackage{listings}
\usepackage{hyperref}

% --- Colores corporativos (mismos que el resto de docs del proyecto) ---
\definecolor{moss}{HTML}{005F73}
\definecolor{moss2}{HTML}{0A9396}
\definecolor{mosslight}{HTML}{D7E8E6}
\definecolor{clay}{HTML}{BB3E03}
\definecolor{sun}{HTML}{EE9B00}
\definecolor{ink}{HTML}{2B2D31}
\definecolor{codebg}{HTML}{F4F6F6}

\hypersetup{
  colorlinks=true,
  linkcolor=moss,
  urlcolor=moss,
  pdftitle={Guía de Defensa - Papaws},
  pdfauthor={Equipo Papaws}
}

\titleformat{\section}{\Large\bfseries\color{moss}}{\thesection.}{0.6em}{}
\titleformat{\subsection}{\large\bfseries\color{clay}}{\thesubsection}{0.6em}{}

% --- Caja para rutas / código corto en línea ---
\newtcbox{\rutabox}{on line, boxrule=0pt, boxsep=0pt, top=1pt, bottom=1pt,
  left=3pt, right=3pt, arc=2pt, colback=codebg, fontupper=\ttfamily\small}

% --- Caja de nota / idea clave ---
\newtcolorbox{nota}{colback=mosslight, colframe=moss, boxrule=0.6pt,
  arc=2pt, left=6pt, right=6pt, top=4pt, bottom=4pt}

% --- Caja para corregir / precisar (acento naranja) ---
\newtcolorbox{precision}{colback=sun!10, colframe=sun!80!black, boxrule=0.6pt,
  arc=2pt, left=6pt, right=6pt, top=4pt, bottom=4pt}

% --- Pregunta / respuesta / dónde mirar ---
\newcommand{\pregunta}[1]{\par\smallskip\noindent{\color{clay}\bfseries P:\ }{\bfseries #1}\par\nobreak}
\newcommand{\dondemirar}[1]{\par\smallskip\noindent{\footnotesize\itshape\color{moss2}\hangindent=1.2em Dónde mirar: #1}\par}

% Holgura global: las rutas en \rutabox son cajas que no se cortan; esto evita que
% se salgan al margen permitiendo estirar los espacios entre ellas.
\sloppy
\setlength{\emergencystretch}{3em}

% --- Estilo de bloque de comandos de shell ---
\lstdefinestyle{shell}{
  basicstyle=\ttfamily\footnotesize,
  backgroundcolor=\color{codebg},
  frame=single, framerule=0pt,
  breaklines=true, columns=fullflexible,
  xleftmargin=6pt, xrightmargin=6pt,
  aboveskip=6pt, belowskip=6pt,
  literate={á}{{\'a}}1 {é}{{\'e}}1 {í}{{\'i}}1 {ó}{{\'o}}1 {ú}{{\'u}}1 {ñ}{{\~n}}1
}

% --- Portada simple con banda verde ---
\newcommand{\portada}{%
\thispagestyle{empty}
\noindent
\begin{tcolorbox}[colback=moss, colframe=moss, arc=3pt, left=14pt, right=14pt,
  top=18pt, bottom=18pt]
  {\color{white}\Huge\bfseries Guía de Defensa}\\[4pt]
  {\color{white}\LARGE\bfseries Papaws}\\[10pt]
  {\color{mosslight}\large Sistema de microservicios distribuidos (.NET + Cassandra + React)}
\end{tcolorbox}
\vspace{1.2em}
}

\begin{document}

\portada

% ===================================================================
\section*{Cómo usar esta guía}
% ===================================================================
\addcontentsline{toc}{section}{Cómo usar esta guía}

Esta guía está en formato \textbf{pregunta $\rightarrow$ respuesta}. Cada tema tiene
primero una explicación clara del concepto y después las preguntas que te pueden
hacer, con la respuesta lista para decir y el lugar exacto del código donde se ve.
Búscala por sección o por palabra clave.

\smallskip
\noindent\textbf{Estructura de cada bloque:}
\begin{itemize}[leftmargin=1.4em, nosep]
  \item \textbf{P:} la pregunta probable del docente.
  \item \textbf{Respuesta:} lo que conviene responder, explicado en simple.
  \item \textit{Dónde mirar:} archivo y método para mostrarlo en vivo si lo piden.
\end{itemize}

\smallskip
\begin{precision}
\textbf{Precisión importante sobre el despliegue.} Donde antes esta guía hablaba de
«una base de datos por servicio», lo correcto es «un \emph{keyspace} por servicio».
Los dos keyspaces conviven en un \textbf{único clúster Cassandra de dos nodos} con
\texttt{replication\_factor~=~2}: la separación entre servicios es \textbf{lógica}
(cada servicio abre solo su keyspace), no física. Esto se detalla en las secciones
\hyperref[sec:queries]{13} y \hyperref[sec:trampa]{15}, y \emph{refuerza} la defensa
sobre tolerancia a fallos.
\end{precision}

\tableofcontents
\newpage

% ===================================================================
\section{La arquitectura en una frase}
% ===================================================================

Papaws son \textbf{TRES microservicios .NET independientes} más un frontend React
que los consume. Dos servicios tienen \textbf{keyspace propio} de Cassandra; el
tercero (Reportes) \textbf{no tiene datos propios}. Ningún servicio entra al
keyspace de otro: si necesita datos ajenos, los pide por la API REST.

\begin{itemize}[leftmargin=1.4em]
  \item \textbf{Animales} (puerto \rutabox{8080}): rescatistas, organizaciones y
        animales. Es el dueño de esos datos. También hace el \textbf{login} y emite
        el \textbf{JWT}.
  \item \textbf{Consulta} (puerto \rutabox{8081}): veterinarios, servicios,
        productos (inventario) y consultas clínicas.
  \item \textbf{Reportes} (puerto \rutabox{8082}): \textbf{NO} tiene datos propios.
        Es un \textbf{AGREGADOR}: arma los reportes pidiendo datos por HTTP a los
        otros dos.
  \item \textbf{Front} (\rutabox{Pawpaws-Front}, React + Vite): decide a qué servicio
        ir por una \emph{«zona»} (\texttt{animales} / \texttt{consulta}) o por la base
        de reportes.
\end{itemize}

\pregunta{¿Por qué microservicios y no un monolito?}
Para desacoplar dominios: cada servicio se despliega, escala y falla por separado.
El precio a pagar es que no hay JOINs ni transacciones entre servicios; la
consistencia se cuida a mano al escribir (sección~\ref{sec:integridad}). Es la
decisión central a defender.

\pregunta{¿Cómo sabe el front a qué microservicio pegar?}
Cada llamada declara su \emph{«zona»}. La función \texttt{baseFor()} traduce la zona
a la URL del servicio. Los endpoints concretos (qué ruta, qué método) están
centralizados en un solo archivo. Reportes no es una «zona»: va por su propia
constante (\texttt{REPORTES\_BASE}) y una función aparte.
\dondemirar{\rutabox{Pawpaws-Front/src/api/client.ts} (\texttt{baseFor} / \texttt{apiGet} / \texttt{apiPost}) y \rutabox{Pawpaws-Front/src/api/endpoints.ts}}

\pregunta{¿La regla de oro de la arquitectura?}
Un servicio NUNCA lee ni escribe el keyspace de otro. Toda dependencia cruzada pasa
por la API REST del dueño, reenviando el token JWT del usuario.

% ===================================================================
\section{Los tres flujos clave (quién llama a quién)}
% ===================================================================

Estas son las tres preguntas de recorrido «de punta a punta» más probables. Conviene
poder trazarlas: \textbf{front $\rightarrow$ controlador $\rightarrow$ servicio
$\rightarrow$ tabla de Cassandra}.

\subsection{¿Dónde llaman los reportes a los rescatistas? (C1 / C2)}
\textbf{Recorrido:}
\begin{itemize}[leftmargin=1.4em]
  \item \textbf{Front}: en \rutabox{Reportes.tsx} se elige el reporte y se llama
        \texttt{reportesApi.c1\_rescatistaPorId(id)} o
        \texttt{c2\_animalesPorRescatista(id)}. Pegan al servicio de
        \textbf{Reportes}, no a Animales.
  \item \textbf{Reportes}: entra por \rutabox{RescatistasReportesController.cs} y
        llama a \texttt{ReporteService} (\texttt{C1\_RescatistaPorIdAsync} /
        \texttt{C2\_AnimalesPorRescatistaAsync}).
  \item Aquí «los reportes llaman a los rescatistas»: \texttt{ReporteService} no
        tiene base; usa el cliente HTTP \texttt{AnimalesClient.GetRescatistaByIdAsync(id)}
        $\rightarrow$ \rutabox{GET api/rescatistas/\{id\}} contra Animales. Para C2
        además \texttt{GetAnimalesByRescatistaAsync(id)} $\rightarrow$
        \rutabox{GET api/animales/rescatista/\{id\}}.
  \item \textbf{Animales} responde desde \texttt{RescatistasController}
        $\rightarrow$ \texttt{RescatistaService}, leyendo la tabla
        \texttt{rescatistas\_by\_id}.
\end{itemize}
\dondemirar{\rutabox{Pawpaws-Reportes/Services/AnimalesClient.cs} y \rutabox{ReporteService.cs}; \rutabox{Pawpaws-Animales/Services/RescatistaService.cs}}

\pregunta{¿Y la autenticación no se pierde al saltar de servicio?}
No. El cliente de Reportes \textbf{reenvía} el token JWT del usuario (copia el header
\texttt{Authorization}), así Animales sigue exigiendo login y rol. Ese reenvío está
en el método que arma el request del cliente.
\dondemirar{\rutabox{Pawpaws-Reportes/Services/AnimalesClient.cs} (\texttt{CrearRequest} copia \texttt{Authorization})}

\subsection{¿Dónde llaman los rescatistas a los animales? (clic en «Ver ficha»)}
\begin{itemize}[leftmargin=1.4em]
  \item \textbf{Front}: el botón «Ver ficha» en \rutabox{Rescatistas.tsx} navega a
        \rutabox{/rescatistas/:id} (ruta en \texttt{App.tsx}) y la pinta
        \rutabox{RescatistaDetalle.tsx}.
  \item Aquí «el rescatista llama a sus animales»: \texttt{RescatistaDetalle.tsx}
        hace \texttt{animalesApi.porRescatista(id)} $\rightarrow$
        \rutabox{GET /api/animales/rescatista/\{id\}} contra Animales.
  \item \textbf{Animales}: \texttt{AnimalesController.ObtenerPorRescatista}
        $\rightarrow$ \texttt{AnimalService.ObtenerPorRescatistaAsync}, que lee la
        tabla \texttt{animales\_by\_rescatista}.
\end{itemize}
\dondemirar{\rutabox{Pawpaws-Front/src/pages/RescatistaDetalle.tsx}; \rutabox{Pawpaws-Animales/Services/AnimalService.cs}}

\pregunta{¿Por qué existe una tabla \texttt{animales\_by\_rescatista} aparte de \texttt{animales\_by\_id}?}
Está \textbf{desnormalizada} y \textbf{particionada} por \texttt{rescatista\_id}
justo para que «dame todos los animales de este rescatista» sea UNA sola lectura de
una partición, sin JOIN ni scan. El mismo animal se guarda también en
\texttt{animales\_by\_id} para buscarlo por su propio id. Ambas se escriben juntas en
un \texttt{BatchStatement} para que no se desincronicen.
\dondemirar{\rutabox{Pawpaws-Animales/Data/CassandraSchema.cs} y \texttt{GuardarAnimalAsync} en \rutabox{AnimalService.cs}}

\subsection{¿Dónde llaman los animales a sus consultas?}
\begin{itemize}[leftmargin=1.4em]
  \item \textbf{Front}: \rutabox{AnimalDetalle.tsx} hace
        \texttt{consultasApi.porAnimal(id)} $\rightarrow$
        \rutabox{GET /api/consultas/animal/\{id\}} contra el servicio de
        \textbf{Consulta} (las consultas viven en otro microservicio).
  \item \textbf{Consulta}: \texttt{ConsultasController.ObtenerPorAnimal}
        $\rightarrow$ \texttt{ConsultaService.ObtenerPorAnimalAsync}.
  \item \textbf{En Cassandra}: primero lee la tabla índice
        \texttt{consulta\_codigos\_by\_animal} (particionada por \texttt{animal\_id})
        para sacar los códigos; luego trae cada consulta de
        \texttt{consultas\_by\_codigo}. Es el patrón «tabla índice + lectura por
        clave».
\end{itemize}
\dondemirar{\rutabox{Pawpaws-Front/src/pages/AnimalDetalle.tsx}; \rutabox{Pawpaws-Consulta/Services/ConsultaService.cs}}

\pregunta{Si el animal y la consulta están en servicios distintos, ¿cómo se relacionan sin foreign key?}
La consulta guarda el \texttt{animal\_id}. No hay FK que cruce servicios; la relación
se sostiene por ese id y la consistencia se cuida al CREAR la consulta, validando el
animal por HTTP (sección~\ref{sec:integridad}).

% ===================================================================
\section{Modelado en Cassandra (query-first)}
% ===================================================================

Cassandra no hace JOINs ni filtra eficiente por columnas que no sean la clave de
partición. Por eso se modela «una tabla por consulta de lectura» (\emph{query-first})
y el mismo dato se guarda varias veces (\emph{desnormalización}). La clave de
partición define dónde vive el dato y por qué campo se puede leer rápido.

\smallskip
\noindent\textbf{Ejemplos en el sistema:}
\begin{itemize}[leftmargin=1.4em]
  \item \textbf{Animales}: \texttt{animales\_by\_id} (por id) y
        \texttt{animales\_by\_rescatista} (por rescatista).
  \item \textbf{Consultas}: \texttt{consultas\_by\_codigo} (por código);
        \texttt{consulta\_codigos\_by\_animal} y
        \texttt{consulta\_codigos\_by\_veterinario} (índices para listar);
        \texttt{consulta\_servicios\_by\_codigo} y
        \texttt{consulta\_productos\_by\_codigo} (relaciones de la consulta).
\end{itemize}
\dondemirar{\rutabox{Pawpaws-Animales/Data/CassandraSchema.cs} y \rutabox{Pawpaws-Consulta/Data/CassandraSchema.cs}}

\pregunta{¿Y si quiero saber en qué consultas se usó un servicio? (consulta inversa)}
No se puede con \texttt{consulta\_servicios\_by\_codigo}, porque esa tabla está
particionada por código (solo responde «qué servicios tiene ESTA consulta»). Por eso
se agregó una tabla de \textbf{índice inverso}, \texttt{consulta\_codigos\_by\_servicio},
particionada por \texttt{servicio\_id}, que se mantiene sincronizada al crear, editar
y borrar una consulta. La usa la ficha de servicio.
\dondemirar{\texttt{ObtenerPorServicioAsync} y los \texttt{INSERT}/\texttt{DELETE} del índice en \rutabox{Pawpaws-Consulta/Services/ConsultaService.cs}}

\pregunta{¿No es peligroso duplicar datos?}
El riesgo es la desincronización. Se controla escribiendo TODAS las copias en la
misma operación (\texttt{BatchStatement}) para que entren juntas, y manteniendo los
índices en cada alta/baja/edición.

% ===================================================================
\section{Trazabilidad e historial (eventos inmutables)}
% ===================================================================

Para no perder historia, los cambios se guardan como \textbf{eventos inmutables} en
tablas ordenadas por fecha descendente (\emph{clustering}). En vez de pisar un valor,
se agrega un evento nuevo.

\begin{itemize}[leftmargin=1.4em]
  \item \texttt{eventos\_custodia\_by\_animal}: ingreso del animal y cada
        reasignación entre rescatistas.
  \item \texttt{eventos\_organizacion\_by\_rescatista}: alta y cambios de
        organización del rescatista.
  \item \texttt{eventos\_adopcion\_by\_animal}: adopciones y devoluciones.
\end{itemize}
\dondemirar{Tablas y \emph{backfill} en \rutabox{Pawpaws-Animales/Data/CassandraSchema.cs}; registro en \texttt{AnimalService.cs} y \texttt{RescatistaService.cs}; líneas de tiempo en \texttt{AnimalDetalle.tsx} y \texttt{RescatistaDetalle.tsx}}

\pregunta{¿Por qué el evento guarda el nombre y no solo el id?}
Para que el historial siga siendo legible aunque después se dé de baja ese rescatista
u organización. Si guardara solo el id, al borrarse quedaría un id muerto sin
sentido. Es un \emph{«snapshot»} del nombre en el momento del evento.

\pregunta{¿Qué es el backfill?}
Un proceso \textbf{idempotente} que siembra los eventos iniciales para los datos que
ya existían antes de tener el historial, sin duplicar si se corre dos veces.

% ===================================================================
\section{Integridad entre servicios (sin foreign keys)}
\label{sec:integridad}
% ===================================================================

Como cada servicio tiene su keyspace, no hay claves foráneas que crucen servicios.
La consistencia se cuida al \textbf{ESCRIBIR}:
\begin{itemize}[leftmargin=1.4em]
  \item Al crear una consulta, Consulta valida el animal \textbf{llamando por HTTP} a
        Animales (con reintentos). Si el animal no existe, rechaza.
  \item Esa misma validación trae el \textbf{estado} del animal y rechaza agendar si
        está «Adoptado» (un adoptado no recibe consultas hasta que se registre su
        devolución).
\end{itemize}
\dondemirar{\rutabox{Pawpaws-Consulta/Services/AnimalReferenceService.cs} (\texttt{ObtenerEstadoAsync}) usado en \texttt{CrearAsync} de \rutabox{ConsultaService.cs}}

\pregunta{¿Y si el servicio de Animales está caído cuando creo la consulta?}
El cliente reintenta con \emph{backoff}. Si aun así no hay respuesta, la creación
\textbf{FALLA} con un error controlado, en lugar de crear una consulta que apunte a un
animal que no se pudo verificar. Se prefiere fallar a dejar datos inconsistentes.

\pregunta{¿Cómo se borra un animal sin dejar consultas huérfanas?}
El animal se borra \textbf{físicamente} y, en cascada, se borran sus consultas. Si la
baja de la cascada falla, se propaga el error y el animal NO se borra, para no dejar
consultas apuntando a un animal inexistente.
\dondemirar{\texttt{EliminarAsync} en \rutabox{Pawpaws-Animales/Services/AnimalService.cs} (cascada cross-servicio)}

% ===================================================================
\section{Concurrencia: el stock de productos (LWT)}
% ===================================================================

Si dos consultas descuentan stock del mismo producto a la vez, un \texttt{UPDATE}
normal podría pisar el valor del otro (\emph{lost update}). Se resolvió con
\textbf{compare-and-set} usando \textbf{LWT} (transacciones ligeras) de Cassandra:
\rutabox{UPDATE ... IF stock\_disponible = ?}. Si otro proceso cambió el stock
entremedio, el \texttt{IF} falla, se lee el valor real que devuelve Cassandra y se
reintenta.
\dondemirar{\texttt{AjustarStockAsync} en \rutabox{Pawpaws-Consulta/Services/ProductoService.cs}}

\pregunta{¿Por qué no usar solo un UPDATE común?}
Porque dos \texttt{UPDATE} concurrentes no se enteran uno del otro: el último gana y
se pierde un descuento. El LWT agrega una condición (\texttt{IF}) que solo aplica si
el valor sigue siendo el que leí; si cambió, reintento con el valor nuevo. Es control
de concurrencia \textbf{optimista}.

\pregunta{¿Es caro el LWT?}
Sí, más que un \emph{write} normal porque usa consenso (Paxos) entre réplicas. Se
justifica solo donde hay riesgo real de carrera, como el stock. No se usa para todo.

% ===================================================================
\section{Gastos y precios (refugio sin fines de lucro)}
% ===================================================================

El reporte de gastos lo calcula \textbf{el servicio de Consulta}, porque ahí viven los
precios: \texttt{precio\_base} de los servicios y \texttt{costo\_unitario} de los
productos. Por cada consulta suma el costo de sus servicios más el de los productos
consumidos (\texttt{costo\_unitario}~$\times$~cantidad). El front lo agrupa por mes.

\begin{itemize}[leftmargin=1.4em]
  \item \textbf{Endpoint}: \rutabox{GET /api/consultas/gastos}
        (\texttt{ObtenerGastosAsync} en \texttt{ConsultaService.cs}).
  \item Las consultas \textbf{canceladas} se EXCLUYEN: devuelven el stock y no
        representan gasto real.
  \item Los precios son en bolivianos (Bs) y a precio de costo, porque es un refugio
        sin fines de lucro: no hay margen comercial.
\end{itemize}
\dondemirar{\rutabox{Pawpaws-Consulta/Services/ConsultaService.cs} (\texttt{ObtenerGastosAsync}); \rutabox{Pawpaws-Front/src/pages/Gastos.tsx}}

\pregunta{¿Por qué el cálculo de gastos no está en Reportes?}
Porque todos los datos necesarios (precio de servicios y costo de productos) viven en
Consulta. Calcularlo ahí evita que Reportes tenga que pedir y recombinar todo por
HTTP: se hace donde están los datos.

\pregunta{¿Por qué los números se muestran sin punto de miles?}
Decisión de formato: el punto se reserva para decimales. La moneda se muestra como
«Bs» con el locale \texttt{es-BO} y agrupación desactivada, para que no se confunda
\texttt{1.000} (mil) con \texttt{1,000} (uno coma cero).

% ===================================================================
\section{Refugio y cupo de rescatistas}
% ===================================================================

Existe un rescatista interno llamado \textbf{«Refugio»} (oculto de la gestión) que
actúa como destino institucional de los animales. Cada rescatista normal tiene un
\textbf{cupo máximo} de animales; el Refugio no tiene límite.

\begin{itemize}[leftmargin=1.4em]
  \item Al dar de baja un rescatista, sus animales se reasignan: llenan el cupo del
        rescatista elegido y lo que sobra se deriva automáticamente al Refugio.
  \item Al reasignar un animal suelto, se bloquea si el rescatista destino ya está en
        su cupo. Así se reparten los animales del Refugio sin sobrecargar a nadie.
\end{itemize}
\dondemirar{\texttt{ReasignarAnimalesAsync} y la validación de cupo en \texttt{ActualizarAsync} de \rutabox{Pawpaws-Animales/Services/AnimalService.cs}; constante \texttt{CapacidadMaxima} en \rutabox{Models/Rescatista.cs}}

\pregunta{¿Por qué no asignar todos los animales a un solo rescatista al dar de baja?}
Porque un rescatista no puede hacerse cargo de más animales de los que su cupo
permite. El sistema respeta ese límite y manda el excedente al Refugio, desde donde
luego se reparten de a poco a quienes tengan lugar.

% ===================================================================
\section{Seguridad: JWT y roles}
% ===================================================================

Login con \textbf{JWT} y autorización por \textbf{roles} (Administrador, Encargado de
consultas, Encargado de rescatistas). Cada endpoint se protege con
\rutabox{[Authorize(Roles = ...)]}. Cuando un servicio llama a otro, \textbf{reenvía}
el token del usuario, así la autorización se respeta en toda la cadena.

\begin{itemize}[leftmargin=1.4em]
  \item El secreto de firma del JWT y la contraseña NO están en el repositorio: se
        leen de variables de entorno / \emph{user-secrets}, y \texttt{appsettings.json}
        queda en blanco (\emph{fail-fast} si faltan).
  \item La misma clave de firma es idéntica en los tres servicios: Animales firma,
        los demás validan.
\end{itemize}
\dondemirar{Carpetas \rutabox{Security/} de cada servicio; emisión del token en Pawpaws-Animales (\texttt{AuthController} / \texttt{JwtService})}

\pregunta{¿Cómo respeta Reportes los permisos si no tiene los usuarios?}
No necesita la base de usuarios: confía en el JWT firmado. Valida la firma con la
clave compartida y lee los roles del token. Además reenvía ese token a
Animales/Consulta, que vuelven a validar.

% ===================================================================
\section{Reglas de negocio}
% ===================================================================

Ejemplos concretos que conviene tener a mano:
\begin{itemize}[leftmargin=1.4em]
  \item Una consulta no puede nacer «Completada».
  \item Un animal «Adoptado» solo puede pasar a «Devuelto al refugio», no directo a
        «Disponible».
  \item No se permiten productos (ni organizaciones) con nombre repetido entre los
        activos.
  \item La fecha de ingreso de un animal no puede ser futura.
  \item Un rescatista no puede superar su cupo de animales (sección~\ref{sec:integridad}
        y «Refugio y cupo»).
\end{itemize}

\pregunta{¿Borrado físico o lógico?}
Rescatistas, veterinarios, servicios, productos y organizaciones usan \textbf{borrado
lógico} (\texttt{activo = false}) para no romper referencias históricas. Los animales
SÍ se borran físicamente, y al hacerlo se borran en cascada sus consultas.

% ===================================================================
\section{Paginación}
% ===================================================================

Hay paginación en \textbf{dos niveles}. En el \textbf{backend}, los listados aceptan
\rutabox{?pagina=} y \rutabox{?tamano=} y devuelven un objeto con los ítems de esa
página. En el \textbf{front}, además, las listas largas se paginan del lado del
cliente: 10 elementos por página con controles numerados (1, 2, 3\ldots).

\begin{itemize}[leftmargin=1.4em]
  \item \textbf{Backend}: helper \texttt{Paginar()} y un tamaño por defecto; el front
        pide la primera página con tamaño alto.
  \item \textbf{Front}: hook \texttt{usePaginated} y componente \texttt{Pagination},
        aplicados a animales, productos, rescatistas, servicios, veterinarios,
        organizaciones, consultas y gastos.
\end{itemize}
\dondemirar{\texttt{apiList} en \rutabox{Pawpaws-Front/src/api/client.ts}; \rutabox{Pawpaws-Front/src/hooks/usePaginated.ts} y \rutabox{components/Pagination.tsx}; \rutabox{Common/Paginacion.cs} en cada servicio}

\pregunta{¿Por qué paginar también en el front si el backend ya pagina?}
Porque el front trae la lista y la filtra/ordena en memoria para una UX fluida;
paginar ahí evita renderizar listas gigantes de cientos de tarjetas. Son dos
optimizaciones que se complementan.

% ===================================================================
\section{Reportes (el servicio agregador)}
% ===================================================================

Reportes no tiene base propia: cada reporte combina llamadas HTTP a Animales y/o
Consulta. El ejemplo más completo es \textbf{C20} («Rescatistas y animales por
organización»): pide la organización, luego sus rescatistas, y para CADA rescatista
sus animales; el «join» organización $\rightarrow$ rescatistas $\rightarrow$ animales
se arma \textbf{en memoria} en el servicio de Reportes, porque los datos viven en
servicios distintos.

\begin{itemize}[leftmargin=1.4em]
  \item Un reporte útil de rendimiento es el \textbf{ranking de veterinarios por
        cantidad de consultas (C22)}, paginado.
\end{itemize}
\dondemirar{\rutabox{Pawpaws-Reportes/Services/ReporteService.cs} (\texttt{C20\_OrganizacionDetalleAsync}, \texttt{C22\_VeterinariosPorConsultasAsync})}

\pregunta{¿Y los datos huérfanos en los reportes?}
Se resuelven los nombres a prueba de huérfanos: si una consulta apunta a un
veterinario dado de baja se muestra «(dado de baja)», y si ya no existe,
«(eliminado)», nunca un GUID suelto.
\dondemirar{\texttt{ResolverNombresVetAsync} / \texttt{ResolverNombresAnimalAsync} en \rutabox{ReporteService.cs}}

\pregunta{¿Por qué el listado de consultas no trae servicios ni productos?}
Para evitar un problema \textbf{N+1} caro en una lista. El listado es liviano; los
servicios y productos se cargan solo al abrir el detalle de una consulta.
\dondemirar{\texttt{ObtenerPorCodigoAsync} en \rutabox{Pawpaws-Consulta/Services/ConsultaService.cs}}

% ===================================================================
\section{Ver las queries reales en Cassandra}
\label{sec:queries}
% ===================================================================

Cassandra se despliega como un \textbf{único clúster} (\texttt{papaws-cluster}) de
\textbf{dos nodos} en contenedores Docker: \texttt{papaws-cassandra-animales}
(puerto \rutabox{9042}) y \texttt{papaws-cassandra-consulta} (puerto \rutabox{9142}).
Ambos comparten \texttt{cluster\_name} y \texttt{seeds}, así que forman un solo
anillo. Conviven dos keyspaces —\texttt{papaws\_animales} y \texttt{papaws\_consulta}—
y como ambos usan \texttt{replication\_factor = 2}, \textbf{cada keyspace está
replicado en los dos nodos}. La separación entre servicios es \textbf{lógica} (cada
servicio abre solo su keyspace), no física. Cassandra NO loguea cada query por
defecto; para «verlas» se activa \textbf{TRACING}.

\begin{nota}
\textbf{Implicación de la réplica.} Como los dos servicios declaran ambos nodos como
\emph{contact points} y cada keyspace tiene una copia en cada nodo (RF~=~2), la caída
de \textbf{un} nodo no tumba a nadie: el otro nodo responde. Esto se usa en la
respuesta sobre el punto único de falla (sección~\ref{sec:trampa}).
\end{nota}

\subsection{Forma A — Una query puntual (cqlsh)}
\begin{lstlisting}[style=shell]
docker exec -it papaws-cassandra-animales cqlsh
USE papaws_animales;
TRACING ON;
SELECT * FROM animales_by_rescatista WHERE rescatista_id = <uuid>;
\end{lstlisting}

\subsection{Forma B — Capturar TODAS las queries (la que más sirve)}
\begin{lstlisting}[style=shell]
# Activar el tracing al 100%
docker exec papaws-cassandra-animales nodetool settraceprobability 1

# Usar la app (abrir una ficha, generar un reporte) y luego leer las trazas:
docker exec papaws-cassandra-animales cqlsh -e \
  "SELECT started_at, request, parameters FROM system_traces.sessions LIMIT 50;"

# APAGARLO al terminar (al 100% es caro):
docker exec papaws-cassandra-animales nodetool settraceprobability 0
\end{lstlisting}
\noindent La columna \texttt{parameters} trae el CQL real.

\subsection{Forma C — Inspeccionar datos directos (demostrar la desnormalización)}
\begin{lstlisting}[style=shell]
docker exec papaws-cassandra-animales cqlsh -e \
  "SELECT * FROM papaws_animales.animales_by_rescatista LIMIT 20;"

# Ver todas las tablas con su clave de partición y clustering:
docker exec papaws-cassandra-animales cqlsh -e "DESCRIBE KEYSPACE papaws_animales;"
\end{lstlisting}

\subsection{Forma D — Logs}
\begin{lstlisting}[style=shell]
docker logs -f papaws-animales   # requests HTTP que llegan al microservicio .NET
\end{lstlisting}
\noindent Cruzado con la Forma B muestra el camino completo: request HTTP
$\rightarrow$ servicio $\rightarrow$ query a Cassandra.

% ===================================================================
\section{Resumen: quién llama a quién}
% ===================================================================

\begin{itemize}[leftmargin=1.4em]
  \item \textbf{Front $\rightarrow$ Animales}: rescatistas, organizaciones, animales
        y sus historiales.
  \item \textbf{Front $\rightarrow$ Consulta}: veterinarios, servicios, productos,
        consultas y gastos.
  \item \textbf{Front $\rightarrow$ Reportes}: todos los reportes.
  \item \textbf{Consulta $\rightarrow$ Animales} (HTTP): solo para validar el animal y
        su estado al crear una consulta.
  \item \textbf{Reportes $\rightarrow$ Animales} y \textbf{Reportes $\rightarrow$
        Consulta} (HTTP): para armar cada reporte combinando datos.
  \item \textbf{Ningún servicio entra al keyspace de otro}: siempre por la API REST,
        con el token JWT reenviado.
\end{itemize}

% ===================================================================
\section{Preguntas trampa y cómo salir bien}
\label{sec:trampa}
% ===================================================================

\pregunta{¿Esto es realmente distribuido o son tres apps cualquiera?}
Es distribuido: tres servicios desplegables por separado, con \textbf{keyspaces
independientes}, comunicados por red (HTTP/REST), con fallas y consistencia
gestionadas explícitamente (reintentos, validación al escribir, eventos). No
comparten datos lógicos ni memoria.

\pregunta{¿No es contradictorio decir «bases independientes» si comparten clúster?}
No, hay que distinguir dos planos. A nivel \textbf{lógico} cada servicio tiene su
propio keyspace y solo abre ese (separación de datos y de dominio real). A nivel
\textbf{físico} los dos keyspaces viven en el mismo clúster de dos nodos por
simplicidad de despliegue; nada impide ponerlos en clústeres separados sin tocar el
código, porque la frontera real es el keyspace, no el contenedor.

\pregunta{¿Qué pasa si dos servicios se desincronizan (un id que ya no existe)?}
Se asume y se maneja: los reportes resuelven nombres a prueba de huérfanos, y las
escrituras críticas validan antes de crear. La consistencia es \textbf{eventual} y
defendida en cada punto de escritura.

\pregunta{¿Por qué Cassandra y no PostgreSQL?}
Por escala de escritura y disponibilidad: Cassandra reparte por particiones y no
depende de un único nodo. El costo es modelar \emph{query-first} y desnormalizar, que
es justamente lo que mostramos.

\pregunta{¿Dónde está el punto único de falla?}
A nivel de aplicación cada servicio y cada keyspace son independientes. A nivel
físico, los dos keyspaces viven en el mismo clúster Cassandra de dos nodos con
\texttt{replication\_factor = 2}: por eso la caída de \textbf{un} nodo no tumba a
nadie (el otro nodo tiene una copia y ambos servicios tienen los dos nodos como
\emph{contact points}). Si el \emph{servicio} de Animales cae, Consulta no puede crear
consultas nuevas (porque valida el animal) pero el resto sigue; la caída se traduce en
un error controlado, no en datos corruptos.

\pregunta{¿Cómo demuestro en vivo que un servicio no toca el keyspace de otro?}
Mostrando que Reportes no tiene cadena de conexión a ninguna base (solo clientes
HTTP) y que cada keyspace solo lo abre su propio servicio
(\texttt{ConnectAsync(keyspace)} con su propio keyspace en cada \texttt{Program.cs}).
Con la Forma B del tracing se ve que las queries del keyspace
\texttt{papaws\_consulta} solo las origina el servicio de Consulta.

\vspace{1.5em}
\begin{center}\small\itshape
Guía de defensa del proyecto Papaws.\\
Complementaria a «Funcionamiento del proyecto Papaws» y «Esquema de tablas en Cassandra».
\end{center}

\end{document}
